\documentclass[11pt]{article}

\usepackage[margin=1.25in]{geometry}
\usepackage{hyperref}
\usepackage{booktabs}
\usepackage{parskip}
\usepackage{enumitem}
\usepackage{titlesec}
\usepackage{microtype}
\usepackage[T1]{fontenc}
\usepackage[utf8]{inputenc}

\hypersetup{
  colorlinks=true,
  linkcolor=blue,
  urlcolor=blue,
  citecolor=blue
}

\titleformat{\section}{\large\bfseries}{}{0em}{}[\titlerule]
\titleformat{\subsection}{\normalsize\bfseries}{}{0em}{}

\setlist[itemize]{leftmargin=1.5em, itemsep=0.25em}

\title{\textbf{Review: Pre-Analysis Plan}\\
\large ``Out of Touch: How Audience Reactions Shape Elite-Public Disconnect''}
\author{}
\date{}

\begin{document}

\maketitle

\section{Executive Summary}

\begin{itemize}
  \item \textbf{Bottom line:} The PAP has unusually strong theoretical grounding, a clean 2$\times$2 design, and transparent analysis specification; the main gaps are a discrepancy between the preregistered index construction method and the analysis code, an inverted primary/robustness label for the two economic perception moderators, underspecified functional forms for income and education, and a missing preregistration platform/timestamp.

  \item \textbf{Fix index construction spec:} The PAP now says outcomes are averaged and standardized to mean 0, SD 1, but the analysis code uses raw sums without standardization---this discrepancy must be resolved before launch so pilot illustrations and full-study results are on the same scale.

  \item \textbf{Clarify RQ1 confirmatory/exploratory status:} RQ1 is labeled exploratory but is given a formal 85\% power target---a power target implies a confirmatory test, and the framing should be made consistent throughout.

  \item \textbf{Specify income and education functional forms:} The PAP adds income and education as moderators but does not specify that income will be log-transformed or that education will be recoded into years---both choices appear in the code but are unspecified in the PAP.

  \item \textbf{Add registration platform and timestamp:} The PAP contains no mention of where or when it will be registered (OSF, AsPredicted, etc.), which is required for the document to function as a true preregistration.

  \item \textbf{Show-stopper:} None; the design supports strong causal claims and no fundamental identification issue is present.
\end{itemize}

\section{To-Do List}

Items are ordered by severity.

\bigskip

\subsection*{1.\enspace Index construction: PAP--code mismatch}

\textbf{Problem.} The PAP's newly added index construction paragraph states that each index is formed as ``a simple unweighted additive average of its component items\ldots\ then standardized to mean 0, SD 1 using the full post-exclusion sample.'' The analysis code, however, builds the indices as raw sums with no division and no standardization (\texttt{disconnect = will\_of\_ppl + good\_corrupt + leader\_represent + public\_officials\_rev + econ\_reality\_rev + system\_stacked}). This creates two problems: (a) the pilot figures embedded in the PAP---which were generated by the code---are on a raw-sum scale, not a standardized scale, meaning the pilot effect sizes reported in the text (e.g., $\hat{\beta} = 0.35$ for H2b) are not in the units the PAP now claims will be used; (b) on the raw-sum scale, the disconnect index (range 6--30) and the epistemic distrust index (range 3--15) are on different scales, making their effect sizes non-comparable without standardization. If standardization is truly preregistered, the code must be updated to implement it, and the pilot figures should either be regenerated on the standardized scale or explicitly noted as being on the raw scale. If raw sums are preferred, the PAP's index construction paragraph should be revised.

\textbf{Evidence.} Index construction paragraph in Section 5: ``standardized to mean 0, SD 1 using the full post-exclusion sample.'' Code lines 104--106: \texttt{df\$disconnect <- df\$will\_of\_ppl + df\$good\_corrupt + \ldots + df\$system\_stacked} and \texttt{df\$epistemic\_distrust <- df\$experts\_claim + df\$trust\_experts\_rev + df\$news\_trust\_rev}. Pilot result in Section 8: ``weakly positive for epistemic distrust ($\hat{\beta}_3 = 0.35$, $p = 0.4$)''---this is on the raw-sum scale.

\textbf{Recommendation.} Decide which approach is canonical: standardized average (as stated in the PAP) or raw sum (as implemented in the code). The standardized approach has the advantage of making the two indices directly comparable and producing interpretable $d$-like effect sizes. If standardization is chosen, update the code to: (1) divide each sum by the number of items to get a mean, (2) standardize to mean 0, SD 1 across the full post-exclusion sample. Then regenerate all pilot figures and update the reported $\hat{\beta}$ values. Alternatively, if raw sums are preferred for simplicity, delete the standardization language from the PAP and instead note that the indices are on different scales and direct comparison of effect sizes across the two outcomes is inappropriate without conversion. Either way, note in the PAP exactly what the scale of the outcome is so readers can interpret the effect sizes.

\textbf{Minimal addition requested.} Add one sentence to the PAP explicitly specifying the outcome scale and confirm the code matches. If standardization is the approach, add three lines of code (\texttt{scale()} on each index) and regenerate the pilot figures. Report at least one pilot effect size in SD units so reviewers can verify the scale.

\bigskip

\subsection*{2.\enspace Primary moderator labeling inverted in code}

\textbf{Problem.} The PAP designates \texttt{econ\_pre} as the primary moderator and \texttt{econ\_identity} as the robustness check throughout---this is explicitly stated in the Measures section (``econ\_pre; primary moderator'') and in all analysis sections (H1, RQ1, Alternative moderators robustness check). However, the analysis code header states the opposite: ``Primary moderator in the pilot is econ\_pre; the full study replaces this with econ\_identity (not collected in Pilot 4) as primary, with econ\_pre as robustness.'' The econ\_identity stub section also says ``econ\_identity is the primary moderator in the preregistered full study.'' If the code is run as-is for the full study, it will implement the analysis with the moderator roles reversed relative to what the PAP preregisters. This is a labeling error in the code (the PAP was recently updated to make econ\_pre primary), not a design problem, but it could cause confusion when the full study launches and could lead to a violation of the preregistration if not caught.

\textbf{Evidence.} PAP Section 5, Measures: ``National economic perceptions (\texttt{econ\_pre}; primary moderator)'' and ``Economic identity (\texttt{econ\_identity}; robustness moderator).'' Analysis code header (lines 9--10): ``the full study replaces this with econ\_identity\ldots\ as primary, with econ\_pre as robustness.'' Code econ\_identity stub section: ``econ\_identity is the primary moderator in the preregistered full study.''

\textbf{Recommendation.} Update the code header and econ\_identity stub section to reflect the PAP: \texttt{econ\_pre} is the primary moderator; \texttt{econ\_identity} is the robustness check. In the stub section, revise the text to say ``All H1, RQ1, education, and income analyses are replicated with \texttt{econ\_identity\_z} substituted for \texttt{econ\_pre\_z} as a secondary robustness check.'' This is a quick text fix in the code and prevents a launch error.

\textbf{Minimal addition requested.} Update lines 9--10 of the code header and the econ\_identity stub description to match the PAP. No analysis changes required.

\bigskip

\subsection*{3.\enspace Missing preregistration platform and timestamp}

\textbf{Problem.} The PAP contains no mention of where it will be registered, when, or under what identifier. A pre-analysis plan that is not submitted to a time-stamped registry (OSF, AsPredicted, EGAP, ClinicalTrials.gov, etc.) cannot be verified as pre-study, which is the entire value of preregistration. Reviewers and editors cannot confirm that the hypotheses and analysis plan were specified before data collection. The PAP also does not state whether data collection has begun, is ongoing, or has not started, which is relevant for assessing the plausibility of the preregistration.

\textbf{Evidence.} The YAML header shows a date of 2026-02-23 and the document ends at Appendix F with no registration link or identifier. No section addresses where or when the PAP will be made publicly accessible.

\textbf{Recommendation.} Register the PAP on OSF (\href{https://osf.io}{osf.io}) before launching data collection and add the registration URL, registration date, and OSF identifier (e.g., ``osf.io/XXXXX'') to the PAP's title page or introduction. Also add a statement of the data collection status at time of registration (e.g., ``Data collection has not begun as of [date]''). If a Registered Report pathway is being pursued at a journal, note that instead. For pilot data: clarify whether pilot data has already been collected and analyzed (it clearly has, given the pilot results sections), and note that the preregistration covers only the full study.

\textbf{Minimal addition requested.} Add a YAML field or title-page note with the following information: registration platform, registration date, registration URL or identifier, and a statement that the full-study data collection has not begun as of registration. This can be a single sentence.

\bigskip

\subsection*{4.\enspace Income and education functional forms unspecified}

\textbf{Problem.} The PAP recently added income as a robustness moderator in the Alternative moderators section, and education appears in the Exploratory analyses section. Neither section specifies how these variables will be operationalized. The analysis code uses \texttt{log(family income)} and education recoded to years of schooling (8 for no qualifications through 20 for PhD), both standardized to $z$-scores. These are non-trivial coding choices: log income versus linear income produces meaningfully different effect size estimates and marginal effects at the tails; the years-of-education mapping is one of several plausible conversions. If these are not preregistered, a reviewer could reasonably ask whether the functional form was chosen after observing results.

\textbf{Evidence.} PAP Robustness Checks: ``econ\_identity and income substituted for econ\_pre as the primary moderator.'' PAP Exploratory Analyses: ``Mismatch effect by education.'' Code lines 77 (\texttt{faminc\_num\_log = log(faminc\_num)}), 92--100 (years mapping), and 117--119 (\texttt{scale()} calls).

\textbf{Recommendation.} Add one sentence to each relevant section specifying: (a) income will be measured as log-transformed family income in dollars, standardized to mean 0, SD 1; missing income is median-imputed with a missingness indicator included as a covariate; (b) education will be recoded to estimated years of schooling using the mapping [specify the categories and years], standardized to mean 0, SD 1. These are the operationalizations already implemented in the code---they just need to be stated in the PAP to close the gap. Specifying median imputation for income is particularly important, as imputation strategy can affect results.

\textbf{Minimal addition requested.} Two sentences in the Robustness/Exploratory sections specifying the functional form, standardization, and income imputation strategy for each variable. The exact mapping for education years can be placed in a footnote.

\bigskip

\subsection*{5.\enspace RQ1 is labeled exploratory but given a confirmatory power target}

\textbf{Problem.} The PAP describes RQ1 (Mismatch Heterogeneity by Economic Satisfaction) as a ``research question'' and notes it is ``exploratory'' in the pilot results table. However, Section 7 (Power Analysis) gives RQ1 a formal power target: ``Approximately 85\% power to detect an interaction of $d = 0.2$ per SD of econ\_pre.'' Power targets are properties of confirmatory tests---they specify the sample size needed to detect a given effect with a given probability if the null is false. Assigning a power target to an exploratory analysis is internally contradictory: exploratory analyses are not powered for specific effects because no inference is being drawn from them. This inconsistency could mislead readers into thinking H3-style predictions are confirmatory, and it weakens the argument that H1 and H2 are the primary confirmatory tests.

\textbf{Evidence.} PAP Section 3: ``RQ1 (Mismatch Heterogeneity By Economic Satisfaction)'' stated as a research question, not a hypothesis. Section 8 pilot results table: labeled as ``RQ1: Mismatch heterogeneity.'' Section 7 power analysis: ``RQ1---Mismatch heterogeneity by economic satisfaction (exploratory): Approximately 85\% power to detect an interaction of $d = 0.2$.''

\textbf{Recommendation.} Either (a) promote RQ1 to a directional hypothesis (H3)---if learning vs.\ activation is a testable prediction with a specific direction, state it as such and commit to a confirmatory test; or (b) remove the formal power target from the RQ1 row in Section 7 and replace with a statement of the minimum detectable effect at the given sample size, labeled as a ``sensitivity analysis'' rather than a power target. Option (b) is lower risk given the mixed pilot results. The MDE paragraph below the power table already contains language appropriate for this (``Null results for these exploratory analyses should be interpreted as consistent with effects smaller than $d \approx 0.35$'')---move the RQ1 power statement to match that framing.

\textbf{Minimal addition requested.} In Section 7, replace the RQ1 ``85\% power'' claim with a minimum detectable effect statement at 85\% power, and label it as a sensitivity calculation rather than a power target. Add a note explaining the distinction between the H1/H2 power targets (confirmatory) and the RQ1 MDE (exploratory sensitivity).

\bigskip

\subsection*{6.\enspace Comment count discrepancy: ``4'' vs.\ ``5 selected''}

\textbf{Problem.} Section 4 (Stimuli) states that ``participants will see 4 user comments (either all negative or all neutral).'' Appendix B then lists ``Negative comments (5 selected)'' and ``Neutral comments (5 selected)'' with five items under each header. Either the stimulus uses 4 of the 5 selected comments (and one was selected but not used), or the stimulus actually shows all 5 and the design description is wrong. This is a small discrepancy but matters for preregistration precision: the exact stimulus is part of what is being preregistered, and a reviewer or replicator needs to know how many comments were shown.

\textbf{Evidence.} Section 4, Stimuli: ``participants will see 4 user comments.'' Appendix B: ``Negative comments (5 selected):'' followed by five numbered items; ``Neutral comments (5 selected):'' followed by five numbered items.

\textbf{Recommendation.} Correct the discrepancy by confirming the actual number shown in the study survey. If participants see all 5, change ``4 user comments'' to ``5 user comments'' in Section 4. If only 4 are shown, remove one item from each list in Appendix B, or note which 4 of the 5 selected are included in the stimulus. Checking the actual Qualtrics survey file (referenced as \texttt{experiment/populism\_and\_news\_-\_pilot4.qsf}) would resolve this immediately.

\textbf{Minimal addition requested.} A one-word correction to the comment count in Section 4 and/or clarification in Appendix B about which comments are actually shown to participants.

\bigskip

\subsection*{7.\enspace Block randomization on econ\_pre not addressed in analysis}

\textbf{Problem.} Section 4 states that ``participants will further be block-randomized based on their responses to the economic perceptions question (\texttt{econ\_pre}).'' Block randomization is used to ensure balanced assignment across economic-perception strata, which is valuable for the het-effects analyses. However, the analysis plan does not mention how block assignment will be handled in estimation. Standard practice in blocked experiments is either to include block fixed effects in all models or to use randomization inference that respects the blocking structure. Neither is mentioned. If blocks are small (which they may be, with 5 levels of \texttt{econ\_pre}), including block fixed effects may be important for unbiased inference. Ignoring the blocking structure means the analysis is valid but potentially inefficient.

\textbf{Evidence.} Section 4: ``Participants will further be block-randomized based on their responses to the economic perceptions question (\texttt{econ\_pre}).'' Analysis plan Section 6: no mention of block fixed effects or blocking-aware inference. The analysis code's control specification does not include a block indicator.

\textbf{Recommendation.} Add a note to Section 6 specifying how block assignment will be handled. The simplest approach is to include \texttt{econ\_pre} as a covariate in all models---which it already is in Specification 1---and note that this implicitly absorbs block-level mean differences. A more rigorous approach would include block fixed effects (one indicator per \texttt{econ\_pre} stratum). Since \texttt{econ\_pre} is a 5-point scale, this adds only 4 binary indicators. The PAP should state explicitly: ``Because randomization is blocked on \texttt{econ\_pre}, \texttt{econ\_pre} is included as a covariate in all primary specifications, which is sufficient to recover the SATE under block randomization (Gerber and Green 2012).'' This one sentence closes the gap.

\textbf{Minimal addition requested.} One sentence in Section 6 explaining how the blocking is handled in estimation and citing an appropriate reference. No change to the analysis code is needed.

\bigskip

\subsection*{8.\enspace Missing Specification 2, Specification 3, IV/CACE, and stimulus-specific robustness implementations}

\textbf{Problem.} The PAP preregisters four robustness checks beyond the main analyses: (1) Specification 2 (saturated controls), (2) Specification 3 (unadjusted), (3) IV/CACE for H1, and (4) stimulus-specific effects by article variant. The analysis code implements only Specification 1 and contains stub or no sections for the others. While the code is labeled as a pilot illustration, a reviewer who reads the PAP alongside the code will notice these gaps. For the full study, all four must be implemented to honor the preregistration. Additionally, the H1 full-sample replication and H2 within-arm check are preregistered but absent from the code.

\textbf{Evidence.} PAP Robustness Checks section lists Specification alternatives, IV/CACE for H1, Stimulus-specific effects, H1 full-sample replication, and H2 within-arm check. Code: only one \texttt{feols()} call per hypothesis, no IV estimation, no by-stimulus loop, no unadjusted or saturated models.

\textbf{Recommendation.} Before the full study launch, add code stubs for all preregistered robustness checks. For Spec 2 and 3, this is trivial: run the same \texttt{feols()} calls without the controls vector and with the expanded controls vector. For IV/CACE, use an IV estimator (e.g., \texttt{feols()} with \texttt{| endogenous \textasciitilde\ instrument} syntax in fixest, or \texttt{ivreg} from the AER package) where article assignment instruments for \texttt{so\_media} or \texttt{so\_elite\_avg}. For stimulus-specific effects, add a by-article-stimulus loop. Adding these stubs before launch---even without data---ensures the full-study analyst follows the preregistration exactly.

\textbf{Minimal addition requested.} Code stubs (even empty function calls with comments) for each of the five missing robustness analyses. No output is needed at this stage; the stubs just need to exist so they are implemented in the full study.

\bigskip

\subsection*{9.\enspace No preregistered differential attrition analysis}

\textbf{Problem.} The PAP preregisters exclusion criteria (attention checks, speeder cutoff) and notes that ``exclusions are applied before all analyses; excluded respondents are described in Appendix C.'' However, it does not preregister any test for differential attrition---whether exclusion rates differ systematically across treatment conditions. Differential attrition is a threat to internal validity even in randomized experiments: if the positive-article arm produces more excludees (e.g., because it generates more frustration), the remaining sample in that arm is a non-random subset. The PAP should preregister a check (e.g., a regression of exclusion indicator on treatment assignment) and specify what action will be taken if differential attrition is detected (e.g., Lee bounds).

\textbf{Evidence.} Section 4, Exclusion criteria: describes two attention checks and a 180-second cutoff. No mention of differential attrition testing in Section 6 (Analysis Plan) or Robustness Checks.

\textbf{Recommendation.} Add a brief paragraph to Section 6 or the Robustness Checks section stating: (a) attrition rates will be compared across the four cells using a regression of an exclusion indicator on treatment assignment; (b) if differential attrition exceeds a threshold (e.g., a 5-percentage-point difference between any two cells), Lee bounds will be computed for all primary estimates. Lee bounds are a simple, widely accepted method that give worst-case bounds on treatment effects under monotone attrition. Computing them in R requires only a few lines using a package like the \texttt{leebounds} package. The pilot data in Appendix C (7 exclusions out of approximately 300) suggests attrition is low, but it should still be checked systematically.

\textbf{Minimal addition requested.} A single sentence in the Analysis Plan preregistering a differential attrition check (regression of exclusion on treatment) and specifying the threshold that would trigger computation of Lee bounds.

\bigskip

\subsection*{10.\enspace Quota sampling: weighting strategy unspecified}

\textbf{Problem.} The PAP plans to quota sample to achieve even thirds on party ID and an even split on education. Quota sampling is described as ``correcting'' Prolific's demographic skew, but the analysis plan does not specify whether survey weights will be applied and, if so, how they will be constructed. If the achieved sample exactly matches the quotas (which is the goal), then no weights are needed. But in practice, even with quotas, some imbalance will remain. The PAP should state explicitly whether analyses will be unweighted (relying on quotas to achieve approximate balance) or whether post-stratification weights will be applied, and if the latter, specify the target population.

\textbf{Evidence.} Section 4, Quota-Sampling: ``I plan to quota sample to meet even third distribution on party ID\ldots\ and even split on college degree.'' Section 6, Analysis Plan: no mention of weights or post-stratification.

\textbf{Recommendation.} Add one sentence to Section 6 clarifying the weighting approach. The simplest defensible choice is: ``No survey weights are applied; quota sampling is used to achieve approximate balance on party ID and education by design, and both variables are included as covariates in all preregistered specifications.'' This closes the gap without adding complexity. If weights are desired for external validity, specify the target marginals (e.g., Current Population Survey) and the weighting method (raking, propensity-score trimming, etc.) before data collection.

\textbf{Minimal addition requested.} One sentence in Section 6 specifying the weighting approach (or explicit statement that analyses are unweighted).

\section{Minor Issues}

These items do not require full discussion but should be addressed before submission.

\begin{itemize}
  \item \textbf{H1 equation redundancy in code:} The H1 model is estimated within the neutral-comments arm (\texttt{df.neu}) but includes \texttt{comment\_neg\_tone} as a predictor, where it is a constant (all \texttt{FALSE}). This variable will be automatically dropped; it should be removed from the formula to avoid confusion.

  \item \textbf{Figure path inconsistency:} The PAP references some figures with \texttt{./figures/} and others with \texttt{figures/} (without the leading \texttt{./}). Both work in Rmarkdown but the inconsistency is worth standardizing.

  \item \textbf{Pilot figure scale note:} The pilot result figures embedded in the PAP (e.g., the H2 mismatch plot, the marginal effects plots) were generated on the raw-sum scale. If the full study uses standardized outcomes, a note should clarify that the pilot figures are on a different scale than the primary study outcomes.

  \item \textbf{``We'' vs.\ ``I'' consistency:} Two remaining instances of ``we'' appear in the document: ``We accept this cost'' in the Design Note on neutral vs.\ no-comments, and ``which does not describe our hypothesis structure'' in the Multiple Comparisons section. Both should be ``I/my'' for consistency with the first-person singular voice used throughout.

  \item \textbf{Cao et al.\ 2026 and Kowalski 2025:} Both citations use years that are at or beyond the current date (early 2026). Confirm these are published, in press, or preprints, and add ``(preprint)'' or ``(forthcoming)'' notation if they have not yet appeared in a peer-reviewed venue. A reviewer will flag unverifiable citations.
\end{itemize}

\end{document}
